%% Generated by md2acm.py from R1_draft.md -- edit the Markdown, then regenerate.
\documentclass[acmsmall,screen,review,anonymous]{acmart}
\usepackage{tabularx}
\usepackage{amssymb}
\newcommand{\inferred}{\textsuperscript{\dag}}
\newcommand{\artifact}[1]{{\footnotesize [artifact: #1]}}
\setcopyright{none}
\acmConference[ISSTA '27]{ACM SIGSOFT International Symposium on Software Testing and Analysis}{September 2027}{Singapore}
\begin{document}
\title{Spec or Bug? Predicting Which Side Developers Repair When a REST API and Its Contract Disagree}
\author{Anonymous Author(s)}
\begin{abstract}
When a REST implementation and its OpenAPI contract disagree, a failing test does not say which artifact is wrong, and the schema-based testing tools we reviewed report the violation without saying which side to repair. We ask which side developers actually repair, and whether that can be predicted.

We mined 248 divergence items from 174 real contract-alignment fixes in 40 open-source Java, Kotlin and Go services. For each item we packaged only the parent state, and had blind LLM coders classify each divergence as \emph{deliberate} (behaviour produced by code written for it) or \emph{accidental} (a side effect of an omission). Confirmatory hypotheses, scorers and packages were hash-frozen before use; exploratory analyses are labelled as such.

In the first confirmatory sample, developers repaired the contract for 26 of 45 items, and divergence kind predicted the repair side: 26/33 deliberate divergences were repaired in the contract, against 0/11 accidental ones (p \textless{} $10^{-4}$). The association held with propositions written by a direction-blind reader (39/48 vs 0/10). It also held in three further pre-registered samples, the last drawn from 15 repositories never used before (33/43 vs 1/18). Pooled over these three samples it was 69/84 vs 6/50 (p = $4.2 \times 10^{-16}$), and it held at the level of fix clusters. Coders from a second model family agreed with the first at $\kappa$ 0.88--1.00.

On real fixes, LLM triage was accurate on average but unstable across samples, with macro-F1 between 0.57 and 0.93 depending on the sample and the labelling method. A pre-registered rule that routes divergences by kind did not beat a forced-choice LLM.

Divergence kind is therefore a robust, interpretable explanation of repair direction. It is not a better triage method than an LLM. Testing tools should stop assuming the implementation is wrong and surface the divergence kind.
\end{abstract}
\maketitle

\section{1. Introduction}
\textbf{The problem.} Schema-based REST testing tools report failures whenever observed behaviour disagrees with the OpenAPI document. Examples are EvoMaster, Schemathesis, RESTler and RESTest. Such a failure has at least two possible origins:
\begin{itemize}
\item the implementation is faulty;
\item the specification is wrong, incomplete, or stale.
\end{itemize}

\textbf{Why existing work does not answer it:}
\begin{itemize}
\item Prior studies classify these failures by hand. Examples are RESTifAI, the TOSEM study of specification/SUT discrepancies, and the RESTest bug reports.
\item Fault-injection work measures how spec faults affect fuzzers, but does not attribute failures. An example is \cite{arXiv2607_12101}.
\item The RESTest authors list automated failure-to-fault classification as an open opportunity.
\end{itemize}

\textbf{Our framing.} We treat the question as \emph{repair-direction triage}: given a divergence, which artifact will the developer change? We answer it with a staged, pre-registered study.

\begin{itemize}
\item \textbf{RQ1 (Visibility).} How often do contract--implementation divergences surface under fuzzing, and is single-sided blame adequate?
\item \textbf{RQ2 (Triage methods).} How well do a zero-shot LLM and an intent-evidence codebook predict the repaired side, on injected faults and on real fixes?
\item \textbf{RQ3 (Predictors).} Does the \emph{kind} of divergence (deliberate vs accidental) predict the repair direction? Does this hold when the divergence is described by someone who does not know the direction?
\item \textbf{RQ4 (Hybrid).} Does routing accidental divergences to the implementation improve LLM triage?
\item \textbf{RQ5 (Robustness).} Does the kind $\rightarrow$ direction association replicate on new samples, with fix capping, with coders from another model family, and on repositories never used before?
\end{itemize}

\textbf{Contributions:}
\begin{enumerate}
\item \textbf{A labelled corpus.} 248 items from real contract-alignment fixes in 40 repositories (E6 25, E7 45, T47 41, T48 40, T49 36, T50 61), labelled with repair direction, plus licence-checked excerpt packages of each item's parent state.
\item \textbf{A blind coding instrument for divergence kind.} Its within-family reliability is $\kappa$ 0.85--1.00, and its cross-family reliability is $\kappa$ 0.88--1.00.
\item \textbf{Confirmatory evidence of the kind $\rightarrow$ direction association.} It holds:
\begin{itemize}
\item under miner-written and direction-blind propositions;
\item in three further pre-registered samples;
\item at the level of fix clusters;
\item on repositories never used before.
\end{itemize}
\item \textbf{A comparison of LLM, codebook and kind-rule triage} on synthetic and real faults. It shows that LLM triage accuracy varies strongly between samples.
\item \textbf{A pre-registered, hash-frozen protocol} for LLM-coded empirical SE studies.
\begin{itemize}
\item Every amendment and deviation is disclosed with its timing.
\item A reproduction script re-runs the frozen E6--E13 scorers and reproduces their outputs byte for byte.
\end{itemize}
\end{enumerate}

\section{2. Background and Related Work}
\emph{Works marked \inferred{} were verified only from their abstracts, web pages or vendor documentation.}

\subsection{The contract as oracle in REST testing}
\begin{itemize}
\item \textbf{Black-box and white-box REST fuzzers use the OpenAPI document as both input model and oracle.} Examples are RESTler ~\cite{P006}, EvoMaster, Schemathesis and RESTest ~\cite{P001}; surveys and tool comparisons, which include EvoMaster, are ~\cite{P002,P003,P004}.
\item \textbf{The usual oracles are 5xx responses and schema violations.} Their findings are often of doubtful severity ~\cite{P015}.
\item \textbf{Spec--response disconformity dominates what such tools report.}
\begin{itemize}
\item In fifteen days of online testing of 13 industrial APIs, 63\% of failures and 50\% of faults were disconformities between the specification and the response ~\cite{P001}.
\item Nearly all of more than 1,000 public specifications showed mismatches ~\cite{P004}.
\item Faulty or under-specified schemas are a main obstacle to fuzzer coverage ~\cite{P003}.
\end{itemize}
\item \textbf{Such a failure has two possible causes:} the implementation is wrong, or the contract is wrong. None of the tools we reviewed predicts which.
\item \textbf{The RESTest authors name the gap.} They list automated classification of failures into faults as an open opportunity (replication package of ~\cite{P001}).
\end{itemize}

\subsection{Contract quality, inference and repair}
\begin{itemize}
\item \textbf{Inference.} Static and path-sensitive analyses infer specifications from code ~\cite{P018,P019}. They find many undocumented endpoints, parameters and responses, and developers confirmed some of the inconsistencies found ~\cite{P018}.
\item \textbf{Repair.} SpecFix repairs documentation-detectable spec defects for LLM-agent wrapping ~\cite{P032}. These are authentication, base URL and undocumented headers, and it works statically only.
\item \textbf{Fault injection.} Mazhar et al. inject faults into OpenAPI documents to measure their effect on black-box fuzzers ~\cite{arXiv2607_12101}. They hold the implementation fixed and do not attribute or repair failures.
\item \textbf{Runtime reconciliation (\inferred{}).} DualMine reconciles spec-derived constraints with runtime executions ~\cite{arXiv2609_27806}. Its output is a constraint set, not a failure-origin label.
\item \textbf{Shared assumption.} All of this work treats the contract as the artifact to be improved. None asks whether the developer would instead change the code.
\end{itemize}

\subsection{Failure diagnosis and triage}
\begin{itemize}
\item \textbf{FDRRestTest (ESEM 2026)} diagnoses failed requests into seven categories and repairs the \emph{test}: value substitution, prerequisite insertion, reordering. It does not attribute a failure to the spec or the implementation.
\item \textbf{Manual labelling elsewhere.} RESTifAI ~\cite{arXiv2512_08706} and LLM test-amplification studies label failing tests by hand as Bug, Enhancement or Invalid. Some "Invalid" cases are business rules missing from the contract, which shows that the distinction matters in practice and is still made manually.
\item \textbf{Practitioner drift tools (\inferred{})} compare live traffic against the OpenAPI document and propose patches. Examples are Beeceptor contract drift and api-drift-detector. A human decides whether the drift is intended.
\item \textbf{Their question is ours.} Is the divergence deliberate, so the contract should follow, or accidental, so the code should be fixed?
\end{itemize}

\subsection{Gap}
We found no work that:
\begin{enumerate}
\item predicts which artifact a developer will repair when contract and implementation disagree;
\item evaluates such a prediction against real developer fixes;
\item identifies an observable property of the divergence that predicts the repair direction.
\end{enumerate}

This paper addresses all three.

\section{3. Study Design}

\subsection{Overview}
\begin{table}[t]
\small
\begin{tabularx}{\linewidth}{lXXX}
\toprule
Stage & Data & Purpose & Status \\
\midrule
Pilot (E1) & scout-api + 3 SUTs; spec mutants and code mutants; fuzzing & RQ1 & descriptive \\
E4 & injected faults (spec and code side) & RQ2, synthetic & confirmatory \\
E5 & code mutants; codebook v3 & RQ2, safety of the codebook & confirmatory \\
E6 & 18 real fixes, 11 repos & RQ2, real & confirmatory \\
E7 & 45 new items from 36 real fixes, 15 repos (+25 E6 items, exploratory) & RQ3 & confirmatory \\
E8 / E8b & E7 items + 41 new (T47) with direction-blind propositions & RQ3, leakage control & confirmatory (pooled) \\
E9 & 85 previously mined items (44 E7 + 41 T47), newly labelled & RQ4 & confirmatory (INSUFFICIENT) \\
E10 & 40 new items (T48), 21 fixes & RQ4 fallback, RQ5 & confirmatory \\
E11 & the 125 earlier confirmatory items, re-coded by a second model family & RQ5 & confirmatory \\
E12 & 36 new items (T49), at most 2 per fix & RQ5 & confirmatory \\
E13 & 61 new items (T50/T50b) from 15 repositories never used before & RQ5 & confirmatory \\
\bottomrule
\end{tabularx}
\end{table}

\subsection{Mining real contract-alignment fixes}
Protocol sources: \texttt{09\_collab/\allowbreak{}tasks/\allowbreak{}T42.md}, \texttt{T45.md}, \texttt{T46.md}, \texttt{T47.md}.
\begin{itemize}
\item \textbf{Repositories:}
\begin{itemize}
\item bare, read-only clones of OSS REST projects (Java, Kotlin, Go);
\item each has a contract surface: a checked-in OpenAPI/Swagger file, or swagger/springdoc/JAX-RS annotations.
\end{itemize}
\item \textbf{Candidates:}
\begin{itemize}
\item \textbf{Commit selection:} commits that touch the contract surface and/or the behaviour surface of the same operation, with a correction cue in the message (e.g. \texttt{swagger}, \texttt{openapi}, \texttt{spec}, \texttt{required}, \texttt{status code}, \texttt{nullable}, \texttt{align}, \texttt{fix}, \texttt{correct}).
\item \textbf{Exclusions:} merges, mass reformatting, dependency bumps and new-endpoint additions.
\end{itemize}
\item \textbf{Ranking:}
\begin{itemize}
\item Candidates are ranked deterministically by cue strength and by whether an operation-specific path is present.
\item The ranking gives no bonus to either direction.
\item The pool is hash-frozen before review.
\item From T46 on, a precise path filter excluded CI, config and workflow YAML files.
\end{itemize}
\item \textbf{Review:}
\begin{itemize}
\item Candidates are reviewed in rank order, with no direction targeting or quotas.
\item Stopping rule: a fixed number of accepted items, or a per-repo review budget.
\end{itemize}
\item \textbf{Acceptance:} the commit resolves a specific, pre-existing mismatch between contract and behaviour for an existing operation.
\begin{itemize}
\item New features with their docs are not accepted.
\item Changes that move both artifacts to new behaviour are not accepted either.
\end{itemize}
\item \textbf{Propositions:}
\begin{itemize}
\item Written from the parent state only, in a fixed neutral template: \texttt{"\textless{}METHOD\textgreater{} \textless{}path\textgreater{}: contract says \textless{}X\textgreater{}; implementation does \textless{}Y\textgreater{}"}.
\item Each proposition has parent-side locators and does not mention which side changed.
\end{itemize}
\item \textbf{Repair direction:}
\begin{itemize}
\item \textbf{SPEC:} only the contract changes;
\item \textbf{IMPL:} only the behaviour changes;
\item \textbf{BOTH:} excluded from the primary analyses.
\end{itemize}
\item \textbf{Miner:} a Codex agent that knew neither hypothesis. It wrote each proposition from the parent state and recorded a direction key.
\item \textbf{Cap:} at most 6 items per repository in T45--T46; from T47 on, a combined cap of 10.
\item \textbf{Later mining (T48--T50):}
\begin{itemize}
\item T48 kept the T46 pool.
\item T49 added a per-fix cap of 2 items.
\item T50 drew a new frozen pool from 20 repositories selected by metadata only. Selection required a checked-in OpenAPI document, $\geq$ 300 commits and an OSI licence. T50 reviewed candidates round-robin across repositories and checked the licence at each item's parent commit.
\end{itemize}
\item \textbf{Upstream code was never executed.}
\end{itemize}

\subsection{Packages and blinding}
\textbf{Item packages:}
\begin{itemize}
\item Each item package contains only the parent state: a contract excerpt and at most 8 implementation files.
\item It has no fix, commit message, date, SHA or history.
\end{itemize}

\textbf{Role separation (all roles are fresh LLM sessions confined to one directory):}
\begin{table}[t]
\small
\begin{tabularx}{\linewidth}{lXXX}
\toprule
Role & Knows the hypothesis? & Knows the direction? & Sees \\
\midrule
Miner & no & yes (writes the key) & repo history \\
Coder (GPT family, C1--C14) & no & no & proposition + package \\
Coder (Claude family, K1--K6) & no & no & proposition + package \\
Proposition-only coder (P1--P4) & no & no & proposition only \\
Direction-blind writer (W1, W2) & no & no & operation + package \\
Matcher (J1, J2) & no & no & W and miner propositions \\
Triage labeller (Y, G, L1--L6, F1--F4) & no & no & proposition + package \\
\bottomrule
\end{tabularx}
\end{table}

\textbf{Protocol controls:}
\begin{itemize}
\item Boundary compliance was checked in every log.
\item Every plan, task text, package tree and scorer was SHA-256 frozen and committed before the run that used it.
\item One frozen scorer received a mechanical fix after it crashed on first load (E12 D1: a CSV field-size limit). The fix changed no logic and came before any output.
\item Every freeze was re-verified before scoring.
\end{itemize}

\subsection{Divergence-kind instrument}
The coders read "the implementation does Y" and judge whether Y comes from project code explicitly written to produce it.

\begin{itemize}
\item \textbf{DELIBERATE:} Y comes from explicit code written for it, such as:
\begin{itemize}
\item a guard, validation or required flag;
\item an explicit status code or response;
\item a serialization, naming or type choice;
\item an explicit branch;
\item a declared route.
\end{itemize}
\item \textbf{ACCIDENTAL:} Y is a side effect of an omission or oversight, such as:
\begin{itemize}
\item an unhandled exception or null dereference;
\item a missing branch;
\item an unset header or content type;
\item an unhandled framework default;
\item an ordering error;
\item a reused variable or cache.
\end{itemize}
\item \textbf{UNCLEAR:} the artifacts do not settle it.
\end{itemize}

\textbf{Requirements:} every code has a parent \texttt{file:line} locator and a rationale of at most 40 words.

\textbf{Secondary field, contract specificity:}
\begin{itemize}
\item \textbf{SPECIFIC:} the contract has an explicit, dedicated statement about this behaviour.
\item \textbf{GENERIC:} the contract's position comes from boilerplate, a default, a shared schema or silence.
\end{itemize}

\subsection{Analyses}
\textbf{Fisher test (H-E7-1, H-E7-2, H-E8-1):}
\begin{itemize}
\item One-sided Fisher exact test on a 2 $\times$ 2 table, with $\alpha$ = 0.05.
\item Wilson 95\% CIs for proportions.
\item Haldane--Woolf CIs for odds ratios.
\item Each arm must contain $\geq$ 8 items before a confirmatory claim is made.
\end{itemize}

\textbf{Triage metrics (RQ2, RQ4):}
\begin{itemize}
\item Macro-F1 over SPEC and IMPL. Any other label, or an abstention, counts as a miss.
\item Decided-accuracy and coverage are also reported.
\end{itemize}

\textbf{Hybrid test (RQ4):} an exact one-sided sign test on discordant items.

\section{4. Results}

\subsection{RQ1: Visibility (pilot, descriptive)}
Source: \texttt{PILOT\_SUMMARY.md} H1, H2, H6, H7, H9.
\begin{itemize}
\item \textbf{Spec faults are often invisible to positive-only fuzzing.} Replay and negative data expose them, but with noise. A mutated operation's signal cleared the noise floor in 34 of 61 fresh runs. For REQUEST\_LOOSE it cleared in only 6 of 25.
\item \textbf{Single-sided blame is inadequate.} 27\% of the decidable baseline defects (30/112; 161 in total) needed both the spec and the code changed.
\item \textbf{The test oracle is itself a source of error:} 16\% of baseline failures (26/161) were caused by the test.
\item \textbf{Code faults surface in fresh fuzzing:} 25 of 31 paired-seed runs, spanning 12 mutants, cleared the noise floor.
\item \textbf{The pattern repeats across fuzzers.} The same fault-class pattern appears under EvoMaster (20/26) and Schemathesis (19/26).
\end{itemize}

\subsection{RQ2: Triage methods}
Sources: \texttt{E4\_RESULTS.md}, \texttt{E5\_*}, \texttt{E6\_RESULTS.md}.

\begin{table}[t]
\small
\begin{tabularx}{\linewidth}{lXX}
\toprule
Setting & M1 zero-shot LLM, macro-F1 & M2 codebook + intent tool, macro-F1 \\
\midrule
Injected faults (E4) & 0.94 & 0.66--0.70 \\
Real fixes (E6, n = 18) & 0.86, 0.88 & 0.72, 0.70 \\
\bottomrule
\end{tabularx}
\end{table}

\begin{itemize}
\item \textbf{The codebook makes no misattributions, but abstains often.} The intent-evidence codebook (E5 v3) made 0 misattributions. It returned UNDECIDABLE whenever independent intent evidence was absent.
\item \textbf{On real fixes the codebook is not safer.} The pre-registered "M2 is safer" criterion (Q3) failed: decided-accuracy was similar, 0.85--0.94.
\item \textbf{Shared errors.} Both methods failed on two items. In each, a lenient contract (an optional field, or a default) met deliberately strict code, and the developer tightened the spec. These items motivated RQ3.
\end{itemize}

\subsection{RQ3: Divergence kind predicts repair direction}

\textbf{E7: miner-written propositions.}
\begin{itemize}
\item \textbf{Data:} 45 new items. By repair direction: SPEC 26, IMPL 18, BOTH 1.
\item \textbf{Reliability:} $\kappa$(kind) = 0.942.
\end{itemize}

\begin{table}[t]
\small
\begin{tabularx}{\linewidth}{lXXXX}
\toprule
 & DELIBERATE & ACCIDENTAL & OR [95\% CI] & p \\
\midrule
P(SPEC) & 26/33 = 0.79 [0.62, 0.89] & 0/11 [0.00, 0.26] & 81.3 [4.3, 1545] & \textless{} $10^{-4}$ \\
\bottomrule
\end{tabularx}
\end{table}

\begin{itemize}
\item \textbf{Sensitivity:} the result holds when restricted to high-confidence mining, when cross-flagged items are excluded, and when only agreed codes are used. It also holds within each popularity stratum and in a leave-one-repo-out analysis (maximum p $6.7 \times 10^{-5}$).
\item \textbf{H-E7-2 (secondary):} P(SPEC | GENERIC) = 15/18 vs P(SPEC | SPECIFIC) = 11/26, p = 0.007.
\end{itemize}

\textbf{Leakage control (A3, proposition-only coder).}
\begin{itemize}
\item A coder that saw only the one-line proposition reproduced the kind code of all 44 primary items (44/44, $\kappa$ = 1.0).
\item The kind is therefore recoverable from the proposition text. That leaves two explanations:
\begin{itemize}
\item kind is an inherent property of the described behaviour;
\item the direction-aware miner phrased the behaviour in a way that leaks the direction.
\end{itemize}
\end{itemize}

\textbf{E8 / E8b: direction-blind propositions.}
\begin{itemize}
\item \textbf{Design:}
\begin{enumerate}
\item A writer who knew neither the direction nor the miner's text wrote its own proposition from the parent package.
\item A matcher kept only the items where the writer described the same divergence as the miner (SAME).
\item Fresh coders then coded the writer's text with the files.
\end{enumerate}
\item \textbf{E8 result:} on SAME items the kind code was unchanged in 35/35. E8 alone was INSUFFICIENT, because its accidental arm had 6 items, below the threshold of 8.
\item \textbf{E8b (pre-registered, one further look, pooled with E8):}
\end{itemize}

\begin{table}[t]
\small
\begin{tabularx}{\linewidth}{lXXXX}
\toprule
 & DELIBERATE & ACCIDENTAL & OR [95\% CI] & p \\
\midrule
P(SPEC), pooled SAME & 39/48 = 0.81 [0.68, 0.90] & 0/10 [0.00, 0.28] & 87.3 [4.7, 1626] & $1.8 \times 10^{-6}$ \\
\bottomrule
\end{tabularx}
\end{table}

\begin{figure}[t]
\centering
\includegraphics[width=\linewidth]{figures/fig_kind_direction.pdf}
\caption{Share of divergences repaired in the contract (SPEC), by divergence kind, with Wilson 95\% CIs. Rows are not independent samples (see text).}
\Description{Forest-style plot of SPEC shares for deliberate and accidental divergences across samples.}
\label{fig:kind}
\end{figure}
\begin{itemize}
\item \textbf{What it shows:} the share of divergences repaired in the contract, by divergence kind, with Wilson 95\% CIs.
\item \textbf{Rows are not independent samples.}
\begin{itemize}
\item E8 + E8b re-describes the E7 items and the T47 items.
\item The T47 row codes the same T47 items from the miner's text.
\item The E6 row is the exploratory set from which the hypothesis came.
\item The E10 row is disjoint from the earlier rows: new items mined by T48, coded by new coders C9 and C10. It is the first sample with accidental divergences repaired in the contract (2/13, from one fix).
\item The two E12 rows are disjoint from all earlier rows: 36 new items mined by T49, at most 2 per fix. The two rows code the same items, once with GPT coders (C11/C12) and once with Claude coders (K3/K4).
\item The two E13 rows are disjoint from all earlier rows: 61 items from 15 repositories never used before, mined by T50. They code the same items with GPT coders (C13/C14) and with Claude coders (K5/K6).
\end{itemize}
\item \textbf{In the E7, E8 + E8b, T47 and exploratory E6 rows, no accidental divergence was repaired in the contract.} In the E10, E12 and E13 rows a few were: 2/13, 3/19 (GPT coders), 1/17 (Claude coders), 1/18 (GPT) and 0/18 (Claude).
\end{itemize}

\begin{itemize}
\item \textbf{The rule's errors are asymmetric.}
\begin{itemize}
\item ACCIDENTAL $\rightarrow$ IMPL had precision 1.00 on the E7--E9 sets, 0.85 in E10, 0.84 in E12 and 0.94 in E13 (GPT kinds).
\item DELIBERATE $\rightarrow$ SPEC was wrong about one time in five before E12. It was right on all 17 items in E12 and on 33/43 (0.77) in E13.
\end{itemize}
\item \textbf{One post-hoc claim was not replicated.} E8 suggested that blind readers miss accidental divergences more often. The pre-registered E8b test gave p = 0.31, so we retract that claim.
\end{itemize}

\subsection{RQ4: Hybrid triage [E9]}
\textbf{Part A (retrospective, exploratory; data seen before):}
\begin{itemize}
\item On the 18 E6 items, M1's errors were all on DELIBERATE items. M1 was correct on 7/7 ACCIDENTAL items and on 7/10 DELIBERATE items, for both labellers.
\item The hybrid therefore changed no labels.
\end{itemize}

\textbf{Part B (confirmatory, 85 new items: SPEC 49, IMPL 36; kinds DELIBERATE 61, ACCIDENTAL 24):}

\begin{table}[t]
\small
\begin{tabularx}{\linewidth}{lXX}
\toprule
Method & Labeller 1 macro-F1 & Labeller 2 macro-F1 \\
\midrule
R (kind rule) & 0.845 & 0.845 \\
M1 (zero-shot LLM) & 0.910 & 0.896 \\
HYB (ACCIDENTAL $\rightarrow$ IMPL, else M1) & 0.910 & 0.903 \\
\bottomrule
\end{tabularx}
\end{table}

\begin{itemize}
\item \textbf{H-E9-1 is INSUFFICIENT.} HYB and M1 disagreed on only 1 item (pooled 1 win, 0 losses), below the pre-registered minimum of 5 discordant items.
\item \textbf{M1 already labels accidental divergences correctly:} 24/24 and 23/24.
\item \textbf{The ACCIDENTAL $\rightarrow$ IMPL rule was correct on all 24 eligible items.} The only prediction it changed was also correct. This is the third sample with no accidental divergence repaired in the spec.
\item \textbf{M1's residual errors are almost all on DELIBERATE items:} 47/61 correct for each labeller. Most misses are UNDECIDABLE on items whose spec was later repaired (12 per labeller).
\item \textbf{Post-hoc (exploratory):} using R only when M1 abstains raised macro-F1 to 0.976 for both labellers, with full coverage. This rule was formed after seeing the data and needs a new pre-registered test.
\end{itemize}

\subsection{RQ4 replication: fallback triage [E10]}
\textbf{Data:} 40 new items (SPEC 21, IMPL 19) from 21 mined fixes in 8 repositories. Final kinds: DELIBERATE 24, ACCIDENTAL 13, UNCLEAR 3 (coder $\kappa$ 0.853). Pre-registered before mining \artifact{\texttt{E10\_PLAN.md}}; a cluster-level sensitivity analysis was frozen before any blind run \artifact{\texttt{E10\_ADDENDUM.md}}.

\begin{table}[t]
\small
\begin{tabularx}{\linewidth}{lX}
\toprule
Method & macro-F1 (labeller pair) \\
\midrule
R (kind rule) & 0.766 \\
M1 (zero-shot LLM, may abstain) & 0.575 / 0.639 \\
FB (M1; R when M1 abstains) & 0.646 / 0.702 \\
M1F (forced choice SPEC/IMPL) & 0.744 / 0.744 \\
\bottomrule
\end{tabularx}
\end{table}

\begin{itemize}
\item \textbf{H-E10-1 (R on M1's abstentions) is INSUFFICIENT.} R was right on 6/6 pairs, but they come from only 3 items in one cluster, below the pre-registered minimum of 8 pairs.
\item \textbf{H-E10-2 (FB better than forced choice) is NOT\_BETTER.} The forced-choice LLM won all 8 discordant pairs (reverse p = 0.0039). By cluster, the difference rests on 2 of 21 clusters (p = 0.25).
\item \textbf{H-E10-3 (kind $\times$ direction) PASSES again:} DELIBERATE $\rightarrow$ SPEC 19/24 vs ACCIDENTAL $\rightarrow$ SPEC 2/13, p = $2.7 \times 10^{-4}$. By cluster: 7/11 vs 1/9, p = 0.025.
\item \textbf{Exploratory: why M1 dropped.}
\begin{itemize}
\item One repository contributed 10 items from 2 fixes, 9 of them repaired in the spec.
\item Every LLM labeller called all 10 IMPL and was wrong on 9; the kind rule was right on 9/10.
\item So M1's errors here were confident wrong labels, not abstentions.
\end{itemize}
\item \textbf{Exploratory: the first accidental divergences repaired in the contract.} Two items from one fix.
\end{itemize}

\subsection{Cross-family replication of kind coding [E11]}
\textbf{Design:} all 125 confirmatory items were re-coded by two blind Claude coders, each started with a fresh context \artifact{\texttt{E11\_PLAN.md}}. Until then every coder had been a GPT-family model.

\begin{table}[t]
\small
\begin{tabularx}{\linewidth}{lX}
\toprule
Result & Value \\
\midrule
Kind $\times$ direction & DELIBERATE $\rightarrow$ SPEC 67/81 vs ACCIDENTAL $\rightarrow$ SPEC 1/39 (Fisher one-sided p = $2.1 \times 10^{-18}$). By cluster: 38/51 vs 1/30, p = $7.1 \times 10^{-11}$ \\
Agreement between the two Claude coders & $\kappa$ = 0.91 \\
Agreement with the GPT final kinds & 116/118, $\kappa$ = 0.96 \\
\bottomrule
\end{tabularx}
\end{table}

\textbf{The association therefore does not depend on the coders' model family.} Both families read the same miner-written propositions, though, so this agreement cannot rule out a shared textual cue. E8b, with direction-blind propositions, addresses that cue. Planned human coding will give a further check.

\subsection{Fix-capped extension with two model families [E12]}
\textbf{Design:}
\begin{itemize}
\item 36 new items (SPEC 20, IMPL 16) from 24 fixes in 6 repositories, mined by a hypothesis-blind miner with at most 2 operations per fix \artifact{\texttt{E12\_PLAN.md}, \texttt{T49\_report.md}}.
\item The pool was exhausted below the 40-item target.
\item Two GPT coders and two Claude coders coded kind, and two GPT labellers (M1 and forced choice) labelled direction.
\end{itemize}

\begin{table}[t]
\small
\begin{tabularx}{\linewidth}{lX}
\toprule
Result & Value \\
\midrule
Kind $\times$ direction, GPT coders & 17/17 vs 3/19, p = $1.3 \times 10^{-7}$. By cluster: 11/11 vs 2/13, p = $3.1 \times 10^{-5}$ \\
Kind $\times$ direction, Claude coders & 17/17 vs 1/17, p = $7.7 \times 10^{-9}$ \\
GPT vs Claude final kinds & 34/34 both-decided items, $\kappa$ = 1.00 \\
R vs forced-choice LLM & 0 discordant items (INSUFFICIENT) \\
Macro-F1 & R 0.917 (GPT kinds) / 0.944 (Claude kinds; coverage 0.94); M1 0.917; M1F 0.917 \\
Pooled E10 + E12 (76 items) & Kind $\times$ direction: 36/41 vs 5/32, p = $3.0 \times 10^{-10}$; clusters 18/22 vs 3/22. Macro-F1: R 0.842, M1 0.751, M1F 0.828 \\
\bottomrule
\end{tabularx}
\end{table}

\textbf{Notes:}
\begin{itemize}
\item \textbf{The association replicated a third time on new items,} with fix capping, by cluster and in both model families \artifact{\texttt{E12\_RESULTS.md}}.
\item \textbf{Exploratory:} R (GPT kinds), M1 and M1F made identical predictions on all 36 items. The E10 advantage of the kind rule over the LLM did not recur, so E12 gives no evidence that kind is a better \emph{triage} signal than a forced-choice LLM.
\item \textbf{R's three errors} were ACCIDENTAL $\rightarrow$ SPEC items from 2 fixes, so accidental $\rightarrow$ contract repair is no longer confined to one fix.
\end{itemize}

\subsection{Generalisation to new repositories [E13]}
\textbf{Design:}
\begin{itemize}
\item \textbf{Repository selection:} 20 repositories none of the earlier experiments had used, chosen by metadata criteria only: a checked-in OpenAPI/Swagger spec, Java, Kotlin or Go, at least 300 commits, and an OSI licence \artifact{\texttt{t50\_repo\_selection/\allowbreak{}SELECTION.md}}.
\item \textbf{Mining:} a hypothesis-blind miner built a frozen pool of 12,832 candidates. It reviewed them round-robin across repositories, with at most 2 items per fix and 10 per repository, and a licence check at each parent commit.
\item \textbf{Sample:} 947 candidates were reviewed. That gave 61 items (SPEC 34, IMPL 27) from 42 fixes in 15 repositories \artifact{\texttt{T50\_report.md}}.
\item \textbf{Plan:} the plan, scorer and repository list were frozen before mining \artifact{\texttt{E13\_PLAN.md}}. The one process deviation was continuing the mining after a premature budget stop; it was decided without seeing directions \artifact{\texttt{E13\_DEVIATIONS.md}}.
\end{itemize}

\begin{table}[t]
\small
\begin{tabularx}{\linewidth}{lX}
\toprule
Result & Value \\
\midrule
Kind $\times$ direction, GPT coders ($\kappa$ 1.00) & 33/43 vs 1/18, p = $2.2 \times 10^{-7}$. By cluster: 22/29 vs 1/13, p = $4.7 \times 10^{-5}$ \\
Kind $\times$ direction, Claude coders ($\kappa$ 0.89) & 34/40 vs 0/18, p = $3.0 \times 10^{-10}$. By cluster: 23/28 vs 0/12, p = $1.1 \times 10^{-6}$ \\
GPT vs Claude final kinds & 55/58, $\kappa$ = 0.88 \\
R vs forced-choice LLM & 0 vs 7 discordant wins; McNemar p = 0.016, favouring the LLM \\
Macro-F1 & R 0.81 (GPT kinds) / 0.86 (Claude kinds; coverage 0.95); M1 0.90; M1F 0.93 \\
Leave one repository out & largest item-level p = $8.7 \times 10^{-6}$ \\
Pooled E10 + E12 + E13 (137 items) & Kind $\times$ direction: 69/84 vs 6/50, p = $4.2 \times 10^{-16}$; clusters 40/51 vs 4/35, p = $4.0 \times 10^{-10}$. Macro-F1: R 0.83, M1 0.82, M1F 0.88 \\
\bottomrule
\end{tabularx}
\end{table}

\textbf{Notes:}
\begin{itemize}
\item \textbf{The association generalises beyond the original repository set.} It held at item and cluster level, in both model families, and without depending on any one repository \artifact{\texttt{E13\_RESULTS.md}}.
\item \textbf{As a triage rule, kind lost to the forced-choice LLM.} All 7 discordant items went to the LLM. Most of R's errors were DELIBERATE items repaired in the code: DELIBERATE $\rightarrow$ SPEC precision was 33/43 = 0.77, against ACCIDENTAL $\rightarrow$ IMPL 17/18 = 0.94.
\end{itemize}

\section{5. Discussion}
\begin{itemize}
\item \textbf{The default assumption is wrong for a large share of fixes.} Tools that assume "the implementation is wrong" would misroute a large share of real fixes. In our sample, developers repaired the contract in more than half of them.
\item \textbf{Accidental divergences can usually be routed to code.} In the first three confirmatory samples no accidental divergence was resolved by changing the contract. E10 found the first two, both from one fix, and E12 found three more, from two fixes (GPT kinds). E13, on new repositories, found one more (1/18). Routing accidental divergences to the code is therefore high-precision, not infallible: pooled E10 + E12 + E13, 6 of 50 went to the contract (precision 0.88 [0.76, 0.94]).
\item \textbf{Deliberate divergences need a human decision.} They lean towards spec repair, but about one in five is still repaired in the code. This matches practitioner tools, where a human decides whether drift is intended.
\item \textbf{How reliable is LLM triage?}
\begin{itemize}
\item In E9 the LLM seemed to recognise the deliberate/accidental distinction implicitly. It almost never needed the accidental override: one label changed. Most of its errors were abstentions on deliberate divergences, with a few confident errors.
\item E10 shows that this picture is sample-dependent. On a new, more clustered sample the LLM confidently mislabelled deliberate divergences as implementation bugs.
\item The kind rule was the more stable of the two: macro-F1 0.85 in E9 and 0.77 in E10, against 0.90 and 0.575--0.639 for the LLM.
\item No single combination rule was best in both samples. Each sample's apparent winner (fallback in E9; kind rule first plus forced choice in E10, macro-F1 0.82) was found post hoc.
\item On the fix-capped E12 sample, the LLM labellers and the kind rule (GPT kinds) agreed on every item (macro-F1 0.917 each).
\item This is consistent with the E10 gap being sample- and cluster-dependent. It gives no evidence of a general advantage for the kind rule.
\item On the new-repository E13 sample, the forced-choice LLM beat the kind rule (0.93 vs 0.81, p = 0.016). Pooled over the three new-item samples, it has the highest macro-F1 (M1F 0.88, R 0.83, M1 0.82).
\item So the kind is best read as an \emph{explanation} of repair direction, not as a triage method that beats LLMs.
\item The practical lesson: report LLM triage accuracy over many fixes and clusters, and use the divergence kind as an interpretable check on it.
\end{itemize}
\item \textbf{Implication for fuzzers.}
\begin{itemize}
\item A runtime failure of the accidental kind (an unhandled exception, a 500, a missing header) is a strong signal of an implementation bug.
\item A clean, explicit rejection that contradicts the spec is more often a stale or under-specified contract.
\end{itemize}
\item \textbf{Contract specificity is a weaker, partly redundant signal.} In E7, divergences against \emph{generic} contract statements were repaired in the contract more often than those against specific ones: 15/18 vs 11/26 (H-E7-2).
\begin{itemize}
\item \textbf{Pooled over the later samples (exploratory, GPT coders):} 37/43 vs 36/79 (p = $7.6 \times 10^{-6}$).
\item \textbf{The effect is unstable across samples.} It was significant in E10 and E12, not in E13, and reversed in E12 with Claude coders.
\item \textbf{It largely overlaps with kind.} Within deliberate divergences it shrinks to 37/40 vs 30/39 (p = 0.052).
\item \textbf{It depends on reading the contract.} The proposition-only control did not reproduce it.
\item So a generic or under-specified contract statement makes contract repair somewhat more likely. Divergence kind remains the primary signal \artifact{\texttt{X\_specificity/\allowbreak{}SPECIFICITY\_EXPLORE.md}}.
\end{itemize}
\end{itemize}

\section{6. Threats to Validity}
\begin{itemize}
\item \textbf{Construct.}
\begin{itemize}
\item The truth label is the developer's repair direction, not an objective fault location.
\item Some SPEC repairs may legitimise a bug.
\end{itemize}
\item \textbf{Same-model coders.}
\begin{itemize}
\item Before E11, the independent divergence-kind coders were all GPT-family models. Pilot labelling and adjudication are excepted. E11, E12 and E13 add Claude coders, who agree with the GPT coders at $\kappa$ 0.96, 1.00 and 0.88. Both families are still LLMs reading the same text, so human validation remains necessary.
\item \textbf{[TODO]} Report the human spot-check: 20 items, coded by the author, with $\kappa$ against the model codes.
\end{itemize}
\item \textbf{Proposition authorship.} Addressed by E8 and E8b. The remaining limit is one writer and one matcher per round, all from the same model family.
\item \textbf{Sampling.}
\begin{itemize}
\item Only correction-cue commits, and only OSS Java, Kotlin and Go projects.
\item Items cluster by repository and fix.
\item The accidental arm is small (n = 10--24 per sample), so the odds-ratio CIs are very wide.
\item E10's 40 items come from 21 fixes, and one repository contributes 10 items from 2 fixes. Cluster-level analyses are reported alongside the item-level ones.
\item E12 capped each fix at 2 items (36 items, 24 fixes), but it exhausted the frozen candidate pool.
\item E13 addressed this with 20 new repositories, chosen mechanically. Only 7\% of its pool was reviewed, and 15 repositories contributed items. The yields therefore say nothing about prevalence.
\end{itemize}
\item \textbf{Hypothesis origin and forking paths.}
\begin{itemize}
\item H-E7-1 came from post-hoc coding of the E6 items, which are reported as exploratory only.
\item All amendments and deviations are disclosed with their timing:
\begin{itemize}
\item E7 A1--A2, before coding;
\item E7 A3, a post-hoc control fixed before it ran;
\item E8 B1, the single further look;
\item E10 addendum, after mining and before any blind run;
\item E12 D1, a mechanical scorer fix before any output;
\item E13 P1, a mining continuation decided without seeing directions.
\end{itemize}
\item Some mining stages previewed candidates beyond their final prefix, as documented in the T46 and T47 reports.
\end{itemize}
\end{itemize}

\section{7. Conclusion}
When an API and its contract disagree, the side developers repair is predictable from what kind of divergence it is.
\begin{itemize}
\item \textbf{Deliberate divergences,} behaviour that the code explicitly produces, are mostly resolved by changing the contract.
\item \textbf{Accidental divergences,} side effects of omissions, are almost always resolved by changing the code.
\item \textbf{How robust this is:} the association passed every pre-registered kind-association test with adequate arms (E7, pooled E8b, E10, E11, E12, E13). It was descriptively consistent in E9, and E8 alone was insufficient. It held:
\begin{itemize}
\item with direction-blind propositions;
\item at the level of fix clusters;
\item with at most two items per fix;
\item on 15 repositories never used before (E13);
\item with coders from a second model family (E11--E13, $\kappa$ 0.88--1.00 across families).
\end{itemize}
\item \textbf{LLM triage is less stable, but it is not worse overall.} A zero-shot LLM can match or beat explicit rules on one sample (E9, E13) and fall well below them on another (E10). Neither pre-registered attempt to combine the LLM with the kind rule beat the LLM alone. Pooled over the new-item samples, a forced-choice LLM had the highest macro-F1.
\end{itemize}

Testing tools should therefore not assume that the implementation is wrong. They should surface the divergence kind as a first-class signal.

\textbf{Future work:}
\begin{itemize}
\item larger human-validated samples;
\item other model families;
\item triage tools evaluated over many fix clusters, not single samples.
\end{itemize}

\section{Data availability}
Upstream licences are documented in \texttt{LICENSE\_REVIEW.md}, checked at each item's parent commit. Excerpts are released under their upstream licences: Apache-2.0, MIT, EPL-2.0, MPL-2.0, GPL-3.0, AGPL-3.0 and CC BY-NC 4.0, with per-directory licence texts for the copyleft ones. One repository later adopted a restrictive licence. Every excerpt from it comes from earlier versions under MIT, Apache-2.0 or CC BY-NC 4.0. As a robustness check, the headline results do not change without the three items closest to the restrictive period: E8b pooled 37/46 vs 0/10, p = $2.6 \times 10^{-6}$.

\textbf{Replication package} (plan: \texttt{REPLICATION\_PACKAGE\_PLAN.md}). It will be released as an anonymous archive for review. It will contain:
\begin{itemize}
\item the plans and freezes;
\item the scorers and keys;
\item all blind codes and labels;
\item the mining scripts;
\item the item excerpts, under per-repository licences.
\end{itemize}

\texttt{reproduce/\allowbreak{}reproduce\_scores.sh} re-runs the frozen E6--E13 scorers and reproduces their 15 committed outputs byte for byte. The pilot, E4/E5 and the exploratory specificity analysis have their own scripts. \texttt{reproduce/\allowbreak{}check\_freezes.py} checks every file hash in the freeze manifests; package-tree hashes were verified separately before each run. The two post-freeze file changes it reports are both disclosed. Transcripts that echo upstream patches or attribution are withheld.

\bibliographystyle{ACM-Reference-Format}
\bibliography{../references}
\end{document}
